\section{A full-support obstruction at rank three}
\label{sec:obstruction}
The following experiment has only two internal memory cuts. The first
input is $x=(x_1,x_2)\in\{-1,1\}^2$, received as one symbol. A command
$j\in\{1,2,3,4,5\}$ is then received, followed by a query $k\in\{1,2\}$.
The machine emits $B\in\{-1,1\}$ in the update receiving $k$. Let $K_1$
count states after $x$ and $K_2$ states after $j$; a retained final answer
has two labels. No further access to $x$ or $j$ is permitted. For
$0<h\le1$ prescribe the conditional mean vectors
\begin{equation}\label{eq:obstruction-array}
\begin{array}{c|ccccc}
j&1&2&3&4&5\\\hline
z_{x,j}&(h,h)&(h,-h)&(-h,0)&(-h/2,hx_1/4)&(-h/2,hx_2/4)
\end{array}.
\end{equation}
The answer law is $\Prb(B=b\mid x,j,k)=(1+b\,z_{x,j}(k))/2$.

The program is clocked and uses arbitrary stochastic rows. Thus every
allowed implementation has common rows $E(s\mid x)$, $T(t\mid s,j)$
and $D(b\mid t,k)$; all past randomness is in $s$ or $t$.

We first record the geometric estimate used by the converse.

\begin{lemma}[A quantitative forced triangle]
\label{lem:forced-triangle}
Put $h=1-\delta$, where $0\le\delta\le1/10$. If a triangle
$P\subset\cube^2$ contains $(h,h)$, $(h,-h)$ and $(-h,0)$, then every
point $(-h/2,y)\in P$ satisfies
\begin{equation}\label{eq:beta}
 |y|\le\beta(\delta):=\frac14+\frac{\delta}{2(1-\delta)}
                      +\frac{\delta}{2-3\delta}<\frac h2.
\end{equation}
\end{lemma}
\begin{proof}
A convex representation of $(h,h)$ includes a vertex $A$ for which
$2-A_1-A_2\le2\delta$, so $A_1,A_2\ge1-2\delta$. A representation
of $(h,-h)$ similarly includes a vertex $B$ with
$B_1\ge1-2\delta$ and $B_2\le-1+2\delta$. These vertices are distinct.
Call the third vertex $D$. In a representation of $(-h,0)$, let $u$
be the total weight on $A,B$. Since their first coordinates are at
least $1-2\delta$ and $D_1\ge-1$,
\[
 -1+\delta\ge -1+2(1-\delta)u,
 \qquad u\le\frac{\delta}{2(1-\delta)}.
\]
The zero second coordinate implies
$|D_2|\le u/(1-u)\le\delta/(2-3\delta)$.
For a point of $P$ with first coordinate $-h/2$, the total weight $v$
on $A,B$ satisfies
\[
 -h/2\ge-1+2hv,
 \qquad v\le\frac{1-h/2}{2h}
               =\frac14+\frac{\delta}{2h}.
\]
Its second coordinate has absolute value at most
$v+(1-v)|D_2|\le\beta(\delta)$. Finally the two fractions in
$\beta$ increase with $\delta$ while $h/2$ decreases, and
\[
 \beta(1/10)=\frac14+\frac1{18}+\frac1{17}
             =\frac{223}{612}<\frac9{20}.
\]
This proves the strict bound throughout the interval, including zero.
\end{proof}

\begin{lemma}[Centered square in an inscribed triangle]
\label{lem:square-trace}
If a triangle in $[-b,b]^2$ contains $[-a,a]^2$, with $a>0$, then
$b\ge2a$.
\end{lemma}
\begin{proof}
Write its three barycentric functions as $\lambda_s(u)=c_s+d_s\cdot u$.
They are nonnegative on $[-a,a]^2$, so $c_s\ge a\norm{d_s}_1$ and
$\sum_sc_s=1$. If the vertices are $v_s$, affine reconstruction gives
$\sum_s v_sd_s^T=I_2$. Taking traces yields
$2\le b\sum_s\norm{d_s}_1\le b/a$.
\end{proof}
This necessary bound is sufficient for our contradiction; no claim of
its attainability in dimension two is used. Its general-dimensional
version belongs to classical cube-absorption geometry; see
Section~\ref{sec:critical}.

\begin{theorem}[Exact incompatible minima and their Pareto frontier]
\label{thm:obstruction}
For every $9/10\le h\le1$, both cuts of \eqref{eq:obstruction-array}
separately admit an exact three-state realization and have ordinary
normalized residual rank three. No realization has $(K_1,K_2)=(3,3)$.
The entire feasible profile region is
\begin{equation}\label{eq:pareto}
 \{(K_1,K_2):K_1\ge3,\ K_2\ge4\}
 \ \cup\ \{(K_1,K_2):K_1\ge4,\ K_2\ge3\}.
\end{equation}
Thus the exact Pareto-minimal profiles are $(3,4)$ and $(4,3)$, and the
minimum full peak is four. For $h<1$ the specified array has full
support. At $h=9/10$ all answer probabilities belong to
$[1/20,19/20]$.
\end{theorem}
\begin{proof}
The four rows at the first cut vary independently through the last two
second coordinates of \eqref{eq:obstruction-array}. They form a
nondegenerate affine square, hence have normalized row rank three. At
the second cut the response vectors lie in $\R^2$ and include three
affinely independent points $(h,h),(h,-h),(-h,0)$. Its normalized row
rank is also three. Factoring through a stochastic register cannot
increase rank, so both state counts are at least three.

Here are matching one-sided constructions. Put $a=h/4$. At the first
cut, three continuation vectors with free coordinate pairs
\begin{equation}\label{eq:parent-triangle}
                (-a,-a),\quad(3a,-a),\quad(-a,3a)
\end{equation}
and the fixed coordinates from \eqref{eq:obstruction-array} enclose the
four actual rows. The triangle is legal because $3a\le3/4$. The
barycentric encoder weights for $(ax_1,ax_2)$ are
\[
 (\lambda_0,\lambda_1,\lambda_2)
   =\left(\frac{2-x_1-x_2}{4},\frac{1+x_1}{4},\frac{1+x_2}{4}\right).
\]
After $j$, represent the resulting two-dimensional mean in the convex
hull of the four corners of $\cube^2$. For a mean $(u,v)$ the weights
on $(\varepsilon_1,\varepsilon_2)$ are
$(1+\varepsilon_1u)(1+\varepsilon_2v)/4$. The decoder returns the
queried corner coordinate. This is a single machine of profile $(3,4)$.

For profile $(4,3)$, first retain $x$. Every actual second-cut mean lies
in the triangle with vertices $h(1,1),h(1,-1),h(-1,0)$. At a mean
$(u,v)$ its weights on those vertices are
\begin{equation}\label{eq:child-triangle-weights}
 \frac{h+u+2v}{4h},\qquad\frac{h+u-2v}{4h},\qquad\frac{h-u}{2h}.
\end{equation}
For the last two columns they are respectively $(1/4,0,3/4)$ or
$(0,1/4,3/4)$; for the first three columns they are the three unit
vectors. These rows followed by the binary decoder prove the second
construction. All tables are rational when $h$ is rational.

Suppose now that a $(3,3)$ machine exists. Let $G_1,G_2,G_3$ be its
full first-cut continuation arrays. Since their convex hull contains
an affine square of dimension two, their affine span equals the affine
span of the four actual rows. Therefore each $G_s$ has the first three
columns of \eqref{eq:obstruction-array} unchanged and has last columns
$(-h/2,u_s)$ and $(-h/2,v_s)$ for some $u_s,v_s$. The first encoder
then forces
\begin{equation}\label{eq:parent-necessary}
 [-a,a]^2\subseteq\conv\{(u_s,v_s):1\le s\le3\}.
\end{equation}
This affine-slice deduction allows arbitrary state continuations; it
neither assumes that they are residuals nor restricts them beforehand
to the residual hull.

The three states at the second cut have decoder mean vectors
$D_1,D_2,D_3\in\cube^2$. Their triangle $P$ contains all actual
second-cut rows, and hence the three points in
Lemma~\ref{lem:forced-triangle}. Every shifted continuation of each
$G_s$ must also lie in $P$, because the update after $j$ is a common
stochastic row into these same three decoder states. In particular,
$(-h/2,u_s),(-h/2,v_s)\in P$. Lemma~\ref{lem:forced-triangle} gives
$|u_s|,|v_s|\le\beta(1-h)<h/2=2a$. This contradicts
Lemma~\ref{lem:square-trace} applied to \eqref{eq:parent-necessary}.
No $(3,3)$ realization exists. Padding the two constructions proves
\eqref{eq:pareto}. The full-support and probability bounds follow
immediately from the displayed means.
\end{proof}

For $h<1$ all residual minimal faces at each cut coincide with the full
terminal-channel face. Thus the support-component rank is just three
at each cut; support separation is not the source of the failure. The
obstruction is simultaneous positivity of the shifted generators.
Theorem~\ref{thm:slice} and Theorem~\ref{thm:obstruction} together locate
the first possible maximum cut rank for this phenomenon at three.
They also distinguish an actual obstruction from a bilinear feasibility
reformulation: the complete feasible profile set is determined here.
